% GENERATED FILE. Edit canonical lessons, then run npm run generate:books.
% canonical-source-hash: fnv1a64:e4ebeb084a84623a
% canonical-lessons: PA-C144-naram, PA-C144-sakht, PA-C144-mushkal, PA-C144-hushiar, PA-C144-alsi

\chapter{Soft, Hard, Difficult, Clever, Lazy}
\label{ch:pa-soft-hard-difficult-clever-lazy}
\hlchaptermodality{144}

\begin{chapteropening}
\textbf{By the end of this chapter:} \emph{I can say soft, hard, difficult, clever and lazy.}

Complete the last lesson of chapter 144: I can say soft, hard, difficult, clever and lazy.
\end{chapteropening}

\section[\texorpdfstring{\pa{ਨਰਮ} (naram)}{ਨਰਮ (naram)}]{\pa{ਨਰਮ} (naram) — soft}

\label{lesson:PA-C144-naram}



\begin{quote}
Before the new one: say the Punjabi for bitter, then the Punjabi for sharp, spicy.
\end{quote}

\subsection*{You'll want to know: \pa{ਨਰਮ}}
\textbf{\pa{ਨਰਮ}} — \emph{naram} — \textquotedblleft{}soft\textquotedblright{}.

\begin{cousinweb}
\textbf{\pa{ਨਰਮ}} (\emph{naram}) is a loan, usually traced to Persian \emph{narm}, \textquotedblleft{}soft\textquotedblright{}. Punjabi breaks up the \emph{-rm} with a vowel.
\end{cousinweb}

\begin{grammarlens}[title={What you've built}]
One more word for this chapter.
\end{grammarlens}

\subsection*{Your turn}
\begin{itemize}
\raggedright
  \item \emph{Say it:} \emph{naram}
  \item \emph{Say it:} \emph{naram}, once more
  \item \emph{Say it:} say \emph{naram}
  \item \emph{Recall:} say the Punjabi for bitter, then the Punjabi for sharp, spicy, then say \emph{naram} again
\end{itemize}

\begin{tcolorbox}[breakable,colback=teal!4,colframe=teal!35!black,title={Before you move on}]
What does \pa{ਨਰਮ} mean? (\textquotedblleft{}Soft\textquotedblright{}.) Say it once more.
\end{tcolorbox}
\section[\texorpdfstring{\pa{ਸਖ਼ਤ} (sakht)}{ਸਖ਼ਤ (sakht)}]{\pa{ਸਖ਼ਤ} (sakht) — hard}

\label{lesson:PA-C144-sakht}



\begin{quote}
Before the new one: say the Punjabi for sharp, spicy, then the Punjabi for soft.
\end{quote}

\subsection*{You'll want to know: \pa{ਸਖ਼ਤ}}
\textbf{\pa{ਸਖ਼ਤ}} — \emph{sakht} — \textquotedblleft{}hard\textquotedblright{}.

\begin{cousinweb}
\textbf{\pa{ਸਖ਼ਤ}} (\emph{sakht}) is a loan, usually traced to Persian \emph{sakht}, \textquotedblleft{}hard, firm\textquotedblright{}.
\end{cousinweb}

\begin{grammarlens}[title={What you've built}]
One more word for this chapter.
\end{grammarlens}

\subsection*{Your turn}
\begin{itemize}
\raggedright
  \item \emph{Say it:} \emph{sakht}
  \item \emph{Say it:} \emph{sakht}, once more
  \item \emph{Say it:} say \emph{sakht}
  \item \emph{Recall:} say the Punjabi for sharp, spicy, then the Punjabi for soft, then say \emph{sakht} again
\end{itemize}

\begin{tcolorbox}[breakable,colback=teal!4,colframe=teal!35!black,title={Before you move on}]
What does \pa{ਸਖ਼ਤ} mean? (\textquotedblleft{}Hard\textquotedblright{}.) Say it once more.
\end{tcolorbox}
\section[\texorpdfstring{\pa{ਮੁਸ਼ਕਲ} (mushkal)}{ਮੁਸ਼ਕਲ (mushkal)}]{\pa{ਮੁਸ਼ਕਲ} (mushkal) — difficult}

\label{lesson:PA-C144-mushkal}



\begin{quote}
Before the new one: say the Punjabi for soft, then the Punjabi for hard.
\end{quote}

\subsection*{You'll want to know: \pa{ਮੁਸ਼ਕਲ}}
\textbf{\pa{ਮੁਸ਼ਕਲ}} — \emph{mushkal} — \textquotedblleft{}difficult\textquotedblright{}.

\begin{cousinweb}
\textbf{\pa{ਮੁਸ਼ਕਲ}} (\emph{mushkal}) is a loan from Arabic \emph{mushkil}, \textquotedblleft{}difficult, doubtful\textquotedblright{}, which reached Punjabi through Persian. Its Arabic root is \emph{sh-k-l} (Lane).
\end{cousinweb}

\begin{grammarlens}[title={What you've built}]
One more word for this chapter.
\end{grammarlens}

\subsection*{Your turn}
\begin{itemize}
\raggedright
  \item \emph{Say it:} \emph{mushkal}
  \item \emph{Say it:} \emph{mushkal}, once more
  \item \emph{Say it:} say \emph{mushkal}
  \item \emph{Recall:} say the Punjabi for soft, then the Punjabi for hard, then say \emph{mushkal} again
\end{itemize}

\begin{tcolorbox}[breakable,colback=teal!4,colframe=teal!35!black,title={Before you move on}]
What does \pa{ਮੁਸ਼ਕਲ} mean? (\textquotedblleft{}Difficult\textquotedblright{}.) Say it once more.
\end{tcolorbox}
\section[\texorpdfstring{\pa{ਹੁਸ਼ਿਆਰ} (hushiār)}{ਹੁਸ਼ਿਆਰ (hushiār)}]{\pa{ਹੁਸ਼ਿਆਰ} (hushiār) — clever}

\label{lesson:PA-C144-hushiar}



\begin{quote}
Before the new one: say the Punjabi for hard, then the Punjabi for difficult.
\end{quote}

\subsection*{You'll want to know: \pa{ਹੁਸ਼ਿਆਰ}}
\textbf{\pa{ਹੁਸ਼ਿਆਰ}} — \emph{hushiār} — \textquotedblleft{}clever\textquotedblright{}.

\begin{cousinweb}
\textbf{\pa{ਹੁਸ਼ਿਆਰ}} (\emph{hushiār}) is a loan, usually traced to Persian \emph{hoshyār}, \textquotedblleft{}alert, wise\textquotedblright{}. It is built on \emph{hosh}, \textquotedblleft{}the senses\textquotedblright{}.
\end{cousinweb}

\begin{grammarlens}[title={What you've built}]
One more word for this chapter.
\end{grammarlens}

\subsection*{Your turn}
\begin{itemize}
\raggedright
  \item \emph{Say it:} \emph{hushiār}
  \item \emph{Say it:} \emph{hushiār}, once more
  \item \emph{Say it:} say \emph{hushiār}
  \item \emph{Recall:} say the Punjabi for hard, then the Punjabi for difficult, then say \emph{hushiār} again
\end{itemize}

\begin{tcolorbox}[breakable,colback=teal!4,colframe=teal!35!black,title={Before you move on}]
What does \pa{ਹੁਸ਼ਿਆਰ} mean? (\textquotedblleft{}Clever\textquotedblright{}.) Say it once more.
\end{tcolorbox}
\section[\texorpdfstring{\pa{ਆਲਸੀ} (ālsī)}{ਆਲਸੀ (ālsī)}]{\pa{ਆਲਸੀ} (ālsī) — lazy}

\label{lesson:PA-C144-alsi}



\begin{quote}
Before the new one: say the Punjabi for difficult, then the Punjabi for clever.
\end{quote}

\subsection*{You'll want to know: \pa{ਆਲਸੀ}}
\textbf{\pa{ਆਲਸੀ}} — \emph{ālsī} — \textquotedblleft{}lazy\textquotedblright{}.

\begin{cousinweb}
\textbf{\pa{ਆਲਸੀ}} (\emph{ālsī}) is a tadbhava, a word inherited from Sanskrit and reshaped by sound change. It comes from \emph{ālasya}, \textquotedblleft{}sloth\textquotedblright{} (CDIAL 1371).
\end{cousinweb}

\begin{grammarlens}[title={What you've built}]
One more word for this chapter.
\end{grammarlens}

\subsection*{Your turn}
\begin{itemize}
\raggedright
  \item \emph{Say it:} \emph{ālsī}
  \item \emph{Say it:} \emph{ālsī}, once more
  \item \emph{Say it:} say \emph{ālsī}
  \item \emph{Recall:} say the Punjabi for difficult, then the Punjabi for clever, then say \emph{ālsī} again
\end{itemize}

\begin{tcolorbox}[breakable,colback=teal!4,colframe=teal!35!black,title={Before you move on}]
What does \pa{ਆਲਸੀ} mean? (\textquotedblleft{}Lazy\textquotedblright{}.) Say it once more.
\end{tcolorbox}
